% Generated draft: requires mathematical and source review.
\documentclass[11pt]{article}
\usepackage[margin=25mm]{geometry}
\usepackage{fontspec}
\setmainfont{Latin Modern Roman}
\usepackage{amsmath,amssymb,mathtools,unicode-math}
\setmathfont{Latin Modern Math}
\usepackage{enumitem,needspace,xurl,hyperref}
\setlength{\parindent}{0pt}
\setlength{\parskip}{5pt}
\setlength{\emergencystretch}{3em}
\newcommand{\N}{\mathbb{N}}
\newcommand{\Z}{\mathbb{Z}}
\newcommand{\Q}{\mathbb{Q}}
\newcommand{\R}{\mathbb{R}}
\newcommand{\C}{\mathbb{C}}
\newcommand{\F}{\mathbb{F}}
\newcommand{\A}{\mathbb{A}}
\newcommand{\PP}{\mathbb{P}}
\newcommand{\E}{\mathbb{E}}
\newcommand{\eps}{\varepsilon}
\newcommand{\dd}{\,\mathrm{d}}
\begin{document}
\section*{UnsolvedMath ID: 801062}
\textbf{Source problem code:} catalogue entry
\subsection{Definitions and mathematical statement}
Write \(0=(0,0)\in\C^2\) for the origin. The supplied record refers to a map \(H_2\) but does not define it. Consequently, the following formalizes the small-cycle property conditionally on specifying the intended map \(H_2:\C^2\to\C^2\).

Using the convention that a cycle is a finite periodic orbit, a cycle of \(H_2\) is a set
\[
\mathcal O(z)=\{z,H_2(z),\ldots,H_2^{n-1}(z)\},
\]
where \(n\geq 1\) is an integer and \(H_2^n(z)=z\). Here \(H_2^0\) is the identity and \(H_2^j\) denotes the \(j\)-fold iterate. Say that \(H_2\) has \emph{small cycles} if every neighborhood \(U\) of \(0\) in the usual topology of \(\C^2\) contains an entire cycle disjoint from \(\{0\}\). Equivalently,
\[
\forall U\subseteq\C^2\text{ a neighborhood of }0,\quad
\exists n\in\Z_{\geq 1}\;\exists z\in\C^2:\quad
H_2^n(z)=z
\quad\text{and}\quad
\forall j\in\{0,\ldots,n-1\},\ H_2^j(z)\in U\setminus\{0\}.
\]
The question is whether the intended \(H_2\) has this property. Its identity and any associated hypotheses must be recovered from the source before this becomes a fully specified mathematical problem. The periodic-orbit convention above allows fixed points other than the origin; the supplied excerpt does not specify a different convention for cycles.
\subsection{Short English statement}
For the intended map \(H_2\), does every neighborhood of the origin in \(\C^2\) contain a whole periodic orbit that avoids the origin? The supplied excerpt omits the definition of \(H_2\), so the question remains conditional on that missing definition.
\subsection{Sources}
\begin{itemize}
\item Supplied local UnsolvedMath registry record, permanent ID \(801062\), exact source problem code \texttt{catalogue entry}, titled ``Henon Maps: Open Problems --- Question 14.'' The record attributes the problem to Berger and Bedford, with proposed year 2023. Its status is supplied as \texttt{open}; this label is not an independently verified mathematical or literature assessment.
\item Berger and Bedford, \emph{Henon Maps: A List of Open Problems} (2023), Question 14, \url{https://arxiv.org/abs/2312.03907}. This reference and attribution are supplied by the registry background. The recorded access date is 2026-10-16 and the extraction method is \texttt{arxiv\_latex}; neither access nor extraction has been independently verified here.
\end{itemize}
\end{document}
